\documentclass[a4paper,11pt]{article}

\usepackage[utf8]{inputenc}
\usepackage[T1]{fontenc}
\usepackage[margin=2.5cm]{geometry}
\usepackage{amsmath}
\usepackage{xcolor}
\usepackage{hyperref}
\usepackage{graphicx}

% code block style
\usepackage{listings}
\usepackage{xcolor}

\definecolor{codegreen}{rgb}{0,0.6,0}
\definecolor{codegray}{rgb}{0.5,0.5,0.5}
\definecolor{codepurple}{rgb}{0.58,0,0.82}
\definecolor{backcolour}{rgb}{0.95,0.95,0.92}

\lstdefinestyle{mystyle}{
    backgroundcolor=\color{backcolour},
    commentstyle=\color{codegreen},
    keywordstyle=\color{magenta},
    numberstyle=\tiny\color{codegray},
    stringstyle=\color{codepurple},
    basicstyle=\ttfamily\footnotesize,
    breakatwhitespace=false,
    breaklines=true,
    captionpos=b,
    keepspaces=true,
    numbers=left,
    numbersep=5pt,
    showspaces=false,
    showstringspaces=false,
    showtabs=false,
    tabsize=2
}
\lstset{style=mystyle}

% Settings
\usepackage[sc]{titlesec}
\newcommand{\comment}[1]{\noindent \textcolor{blue}{\textit{#1}} \par}

% Header information
\title{\Large{\sc{Performance Engineering}}\\[2mm] Lab 6}
\author{Jort Oosterveld S5545781 \& Sergiu Preda S6778542 \\ Group 11}

\begin{document}
\setlength\parindent{0pt}

\maketitle

\section{Data Structure Optimization}

\begin{lstlisting}[language=C++, caption=Data structures Small and Big]
const int n = 16384;

struct Small {
    char flg1;
    double x[7];
    char fig2;
}

struct Big {
    int rank;
    double x[n]
}
\end{lstlisting}

% ================================ a ==================================
\subsection{a.}
\begin{itemize}
    \item \textbf{Data packing:} \\
        \begin{itemize}
            \comment{Assumes that the data structure has default area size per field of 64 bit}
            \item \textbf{Small:} yes, packs the two char flag fields into one 64-bit area, instead of two separate ones. Used in combination with padding here to pad to a 64-bit area. Also used in combination with field reordering to enable the packing, since in the current implementation, the chars cannot be packed when the double array is in between.
            \item \textbf{Big:} no, data packing does not help in this case, since there are no fields that can fit together in the default area that the \texttt{Big} struct allocates per field (64 bit).
        \end{itemize}
    \item \textbf{Data padding:}
        \begin{itemize}
                \item \textbf{Small:} Yes, used in combination with data packing. If used on its own, would increase the size of the struct in memory and would not improve resource usage, but would make memory access more efficient by aligning the fields to a 64-bit size.
                \item \textbf{Big:} yes, used here to align the 32-bit \texttt{rank} to the 64-bit doubles, for more streamlined memory access.
            \end{itemize}
    \item \textbf{Data Field reordering:}
        \begin{itemize}
                \item \textbf{Small:} Yes, used in combination with packing, reordering enables the packing in this case.
                \item \textbf{Big:} no, reordering does not enable optimizations in the \texttt{Big} case.
            \end{itemize}
    \item \textbf{Data Structure Splitting:}
        \begin{itemize}
                \item \textbf{Small:} Yes, splitting flags and array optimizes access to the flags by not loading array into cache each time.
                \item \textbf{Big:} yes, splitting optimizes access to rank in the same way as with small.
            \end{itemize}
\end{itemize}



\section{MPI Point-To-Point Communication Benchmarking}

\end{document}
